%%%%%%%%%%%%%%%%%%%%

% Semester Project for EPFL
% Author: Keng-Yu Chen

%%%%%%%%%%%%%%%%%%%%

%%%%%%%%%%%%%%%%%%%%

% Semester Project for EPFL
% Author: Keng-Yu Chen

%%%%%%%%%%%%%%%%%%%%

\documentclass[a4paper, 12pt]{article}
\usepackage[utf8]{inputenc}
\usepackage[margin=3cm]{geometry}
\setlength{\headheight}{14.49998pt}

% Include Packages
\usepackage{mathrsfs}
\usepackage{color}
\usepackage{amsthm,amsmath,amssymb}
\usepackage{bm}
\usepackage{stmaryrd}
\usepackage{hyperref}
\usepackage{graphicx}
\usepackage{multicol}
\usepackage{wrapfig}
\usepackage{subcaption}
\hypersetup{colorlinks=true,linkcolor=blue,citecolor=blue,filecolor=magenta,urlcolor=blue}

\usepackage{fancyhdr}
\usepackage{thmtools}

\usepackage{algorithm}
\usepackage{algpseudocode}
\renewcommand{\algorithmicrequire}{\textbf{Input:}}
\renewcommand{\algorithmicensure}{\textbf{Output:}}



% Set Self-defined Words
\DeclareMathOperator*{\argmax}{arg\!\max}
\DeclareMathOperator*{\argmin}{arg\!\min}

\newcommand{\negl}{\mathsf{negl}}
\newcommand{\poly}{\mathsf{poly}}
\newcommand{\Adv}{\mathbf{Adv}}

\newcommand{\R}{\mathbb{R}}
\newcommand{\N}{\mathbb{N}}
\newcommand{\Z}{\mathbb{Z}}
\newcommand{\Q}{\mathbb{Q}}

\newcommand\getsdollar{\mathrel{{\leftarrow}\vcenter{\hbox{\tiny\rmfamily\upshape\$}}}}


% Convert block[name] into a theorem-like format
\newtheorem{theorem}{Theorem}

\theoremstyle{definition}
\newtheorem{definition}{Definition}

\theoremstyle{definition}
\newtheorem{assumption}{Assumption}

\theoremstyle{plain}
\newtheorem{corollary}{Corollary}



% TikZ Setting
\usepackage{tikz}
\usetikzlibrary{shapes.geometric, arrows}
\usetikzlibrary{positioning}
\usetikzlibrary{calc}
\usetikzlibrary{backgrounds}

\tikzstyle{st-arrow} = [thick,->,>=stealth]
\tikzstyle{arrow} = [thick,-latex']
\tikzstyle{darrow} = [thick,latex'-latex']

\tikzstyle{gadget} = [rectangle, rounded corners,
	minimum width=2cm,
	minimum height=1cm,
	text centered,
	draw=black
]

\tikzstyle{element} = [rectangle, rounded corners,
	minimum width=2cm,
	minimum height=1.2cm,
	text centered,
	fill=white!70!black,
	draw=black
]
\pgfdeclarelayer{background layer}
\pgfdeclarelayer{foreground layer}
\pgfsetlayers{background layer,main,foreground layer}




% Add Referemces
\usepackage[
	backend=biber,
	style=alphabetic,
	sorting=ynt
]{biblatex}
\addbibresource{reference.bib}


\title{\textbf{Biometrics Authentication: Formalization and Instantiation}}
\author{Keng-Yu Chen}

\date{\today}



\begin{document}

%% Title
\maketitle

%-------------------

%% Header and Foot
\pagestyle{fancy}
\fancyhf{}
\fancyhead[L]{Semester Project}
\fancyhead[C]{Biometric Authentication}
\fancyhead[R]{Keng-Yu Chen}
\fancyfoot[C]{\thepage}

%-------------------

This report formalizes the biometric authentication scheme, including its structure, usage, and security analysis with a security game model.

%-------------------

\section{Preliminaries}
\label{sec:preliminaries}

In this report, we assume

\begin{itemize}
	
	\item $\lambda$ is the security parameter.

	\item $[m]$ denotes the set of integers $\{1, 2, \cdots, m\}$.

	\item $\Z_q$ is the finite field modulo a prime number $q$.

	\item A function $f(n)$ is called \emph{negligible} iff for any integer $c$, $f(n) < \frac{1}{n^c}$ for all sufficiently large $n$. We write it as $f(n) = \negl$, and we may also use $\negl$ to represent an arbitrary negligible function.
	
	\item $\poly$ is the class of polynomial funcions. We may also use $\poly$ to represent an arbitrary polynomial function.
	
	\item We write sampling a value $r$ from a distribution $\mathcal{D}$ as $r \getsdollar \mathcal{D}$. If $S$ is a finite set, then $r \getsdollar S$ means sampling $r$ uniformly from $S$.

	\item The distribution $\mathcal{D}^t$ denotes $t$ identical and independent distributions of $\mathcal{D}$.

	\item A PPT algorithm denotes a probabilistic polynomial time algorithm. Unless otherwise specified, all algorithms run in PPT.

\end{itemize}


\begin{definition}[Functional Hiding Inner Product Functional Encryption]
\label{def:fh-IPFE}
	A \emph{functional hiding inner product functional encryption} (fh-IPFE) scheme $\sf FE$ for a field $\mathbb{F}$ and input length $k$ is composed of PPT algorithms ${\sf FE.Setup}$, ${\sf FE.KeyGen}$, ${\sf FE.Enc}$, and ${\sf FE.Dec}$:

	\begin{itemize}
	
		\item ${\sf FE.Setup}(1^\lambda) \to {\sf msk}, {\sf pp}$: It outputs the public parameter ${\sf pp}$ and the master secret key ${\sf msk}$.
	
		\item ${\sf FE.KeyGen}({\sf msk}, {\sf pp}, \mathbf{x}) \to f_\mathbf{x}$: It generates the functional decryption key $f_\mathbf{x}$ for an input vector $\mathbf{x} \in \mathbb{F}^k$. 
	
		\item ${\sf FE.Enc}({\sf msk}, {\sf pp}, \mathbf{y}) \to \mathbf{c_y}$: It encrypts the input vector $\mathbf{y} \in \mathbb{F}^k$ to the ciphertext $\mathbf{c_y}$. 
	
		\item ${\sf FE.Dec}({\sf pp}, f_\mathbf{x}, \mathbf{c_y}) \to z \in \mathbb{F}$: It outputs a value $z$.
	
	\end{itemize}
	
	\noindent Correctness: The fh-IPFE scheme {\sf FE} is \emph{correct} if $\forall ({\sf msk}, {\sf pp}) \gets {\sf FE.Setup}$, $\forall \mathbf{x}, \mathbf{y} \in \mathbb{F}^k$, we have
	\[
		{\sf FE.Dec}( {\sf pp}, {\sf FE.KeyGen}({\sf msk}, {\sf pp}, \mathbf{x}), {\sf FE.Enc}({\sf msk}, {\sf pp}, \mathbf{y}) ) = \langle \mathbf{x} , \mathbf{y} \rangle \in \mathbb{F}.
	\]

\end{definition}

%-------------------

\section{Formalization}
\label{sec:formalization}


In general, an authentication shceme $\Pi$ associated with a family of biometric distributions $\mathbb{B}$ is composed of the following algorithms.

\begin{itemize}

	\item ${\sf USetup}(1^\lambda) \to {\sf esk}, {\sf psk}, {\sf csk}$: It outputs the enrollment secret key ${\sf esk}$, probe secret key ${\sf psk}$, and compare secret key $\sf csk$.

	\item ${\sf encodeEnroll}^{\mathcal{O}_{\mathcal{B}}}() \to \mathbf{x}$: Given an oracle $\mathcal{O}_{\mathcal{B}}$, which samples biometric data from the distribution $\mathcal{B} \in \mathbb{B}$, it encodes biometric samples as $\mathbf{x}$, the input format for enrollment. 

	\item ${\sf Enroll}({\sf esk}, \mathbf{x}) \to \mathbf{c_x}$: It outputs the enrollment message $\mathbf{c_x}$ from $\mathbf{x}$.

	\item ${\sf encodeProbe}^{\mathcal{O}_{\mathcal{B}}}() \to \mathbf{y}$: Given an oracle $\mathcal{O}_{\mathcal{B}}$, which samples biometric data from the distribution $\mathcal{B} \in \mathbb{B}$, it encodes biometric samples as $\mathbf{y}$, the input format for probe.

	\item ${\sf Probe}({\sf psk}, \mathbf{y}) \to \mathbf{c_y}$: It outputs the probe message $\mathbf{c_y}$ from $\mathbf{y}$.

	\item ${\sf Compare}({\sf csk}, \mathbf{c_x}, \mathbf{c_y)} \to s$: It compares the enrollment message $\mathbf{c_x}$ and probe message $\mathbf{c_y}$ and outputs a score $s$.

	\item ${\sf Verify}(s) \to r \in \{0,1\}$: It is a deterministic algorithm that reads the comparison score $s$ and determines whether this is a successful authentication ($r = 1$) or not ($r = 0$).

\end{itemize}


The usage model we consider is described in Figure \ref{fig:model_enrollment} (on enrollment) and Figure \ref{fig:model_auth} (on authentication).

\begin{itemize}

	\item {\sf User}: The user who enrolls its biometric data and authenticates itself to the server. We assume the user's biometric distribution is $\mathcal{B} \in \mathbb{B}$. 
	\item {\sf Scanner}: A machine to extract the user's biometric data by querying the oracle $\mathcal{O}_{\mathcal{B}}$.
	\item {\sf Device}: A device belonging to the user. In practice, it can be a desktop or a mobile phone. It is responsible for most computation and storing the enrollment secret key {\sf esk} and probe secret key {\sf psk}. It queries $\mathcal{O}_{\mathcal{B}}$ for biometric data through the {\sf Scanner}.
	\item {\sf Server}: The server responsible for authenticating the user. It stores the comparison secret key {\sf csk} and the user's enrollment message $\mathbf{c_x}$. On authentication, it compares the probe message with the registered enrollment message and returns the result.  

\end{itemize}

%\begin{figure}[]
	%\begin{center}
		%\begin{tabular}{l c r c}
			%{\bf User} & 	& {\bf Server} \\

			%\hline

			%\\

			%${\sf esk}, {\sf psk}, {\sf csk} \gets {\sf USetup}(1^\lambda)$ \\

			%$ \mathbf{b} \getsdollar \mathcal{B}, \mathbf{x} \gets {\sf encodeEnroll}^{\mathcal{O}_{\mathcal{B}}}()$ \\

			%$\mathbf{c_x} \gets {\sf Enroll}({\sf esk}, \mathbf{x})$ \\

			%& $\xrightarrow{\quad {\sf csk}, \mathbf{c_x} \quad}$ \\

			%$ \mathbf{b'} \getsdollar \mathcal{B}, \mathbf{y} \gets {\sf encodeProbe}^{\mathcal{O}_{\mathcal{B}}}()$ \\

			%$\mathbf{c_y} \gets {\sf Probe}({\sf psk}, \mathbf{y})$ \\

			%& $\xrightarrow{\quad \mathbf{c_y} \quad}$ \\
			
			%&  & $s \gets {\sf Compare} ({\sf csk}, \mathbf{c_x}, \mathbf{c_y})$ \\

			%& & $r \gets {\sf Verify}(s) $ \\

			%& $\xleftarrow{\quad r \quad}$ & 
		%\end{tabular}
	%\end{center}
	%\caption{Authentication Model with User and Server}
	%\label{table:model}
%\end{figure}

%-------------------

\begin{figure}
\centering

\begin{tikzpicture}[node distance=2cm]

% Gadgets
\node (user) [inner sep=0pt] at (0,8) 
    {\includegraphics[width=0.1\textwidth]{image/user.jpg}};

\node (scanner) [element] at (1.5,6.5) {{\sf Scanner}}; 

\node (device) [element] at (7,6.5) {{\sf Device}};

\node (server) [element] at (12,6.5) {{\sf Server}};

% Texts
% User and Scanner
\draw [arrow] (scanner.west) -| node [align=center, left] {{\sf query} \\ $\mathcal{O}_\mathcal{B}$} (user.south);

\draw [arrow] (user.east) -| node [align=center, right] {{\sf resp.} \\ $\mathbf{b} \gets \mathcal{B}$}  (scanner.north);

% Device
\node [align=center] at ($(device.south) + (0, -1)$) {
    ${\sf esk}, {\sf psk}, {\sf csk} \gets {\sf USetup}(1^\lambda)$ \\
    Store ${\sf esk}, {\sf psk}$        
};
\node [align=center] at ($(device.south) + (0, -2)$) {
	$\mathbf{x} \gets {\sf encodeEnroll}^{\mathcal{O}_{\mathcal{B}}}()$
};
\node [align=center] at ($(device.south) + (0, -3.8)$) {
    $\mathbf{c_x} \gets {\sf Enroll}({\sf esk, \mathbf{x}})$
};

% Scanner and Device
\draw [arrow]  ($(device.west) + (-1.5, -2.6)$) -- node[align=center, above] {{\sf query} \\ $\mathcal{O}_\mathcal{B}$} ($(scanner.east) + (0.5, -2.6)$);
\draw [arrow]  ($(scanner.east) + (0.5, -2.8)$) -- node[align=center, below] {{\sf resp.} \\ $\mathbf{b} \gets \mathcal{B}$} ($(device.west) + (-1.5, -2.8)$);

% Device and Server
\draw [arrow]  ($(device.east) + (1.5, -1.5)$) -- node[align=center, above] {${\sf csk}$} ($(server.west) + (-0.5, -1.5)$);
\draw [arrow]  ($(device.east) + (1.5, -4.5)$) -- node[align=center, above] {$\mathbf{c_x}$} ($(server.west) + (-0.5, -4.5)$);

% Server
\node [align=center] at ($(server.south) + (0, -1)$) {
    Store ${\sf csk}$
};
\node [align=center] at ($(server.south) + (0, -4)$) {
    Store $\mathbf{c_x}$
};


\end{tikzpicture}

\caption{Usage Model on Enrollment}
\label{fig:model_enrollment}
\end{figure}







% Authentication Model
\begin{figure}
\centering

\begin{tikzpicture}[node distance=2cm]

% Gadgets
\node (user) [inner sep=0pt] at (0,8) 
    {\includegraphics[width=0.1\textwidth]{image/user.jpg}};

\node (scanner) [element] at (1.5,6.5) {{\sf Scanner}}; 

\node (device) [element] at (7,6.5) {\sf Device};

\node (server) [element] at (12,6.5) {\sf Server};

% Texts
% User and Scanner
\draw [arrow] (scanner.west) -| node [align=center, left] {{\sf query} \\ $\mathcal{O}_\mathcal{B}$} (user.south);

\draw [arrow] (user.east) -| node [align=center, right] {{\sf resp.} \\ $\mathbf{b} \gets \mathcal{B}$}  (scanner.north);

% Device
\node [align=center] at ($(device.south) + (0, -0.9)$) {
    $\mathbf{y} \gets {\sf encodeProbe}^{\mathcal{O}_{\mathcal{B}}}()$
};
\node [align=center] at ($(device.south) + (0, -2.7)$) {
    $\mathbf{c_y} \gets {\sf Probe}({\sf psk, \mathbf{y}})$
};

% Scanner and Device
\draw [arrow]  ($(device.west) + (-1.3, -1.5)$) -- node[align=center, above] {{\sf query} \\ $\mathcal{O}_\mathcal{B}$} ($(scanner.east) + (0.7, -1.5)$);
\draw [arrow]  ($(scanner.east) + (0.7, -1.7)$) -- node[align=center, below] {{\sf resp.} \\ $\mathbf{b} \gets \mathcal{B}$} ($(device.west) + (-1.3, -1.7)$);

% Device and Server
\draw [arrow]  ($(device.east) + (1, -3.4)$) -- node[align=center, above] {$\mathbf{c_y}$} ($(server.west) + (-1, -3.4)$);
\draw [arrow]  ($(server.west) + (-1, -4.4)$) -- node[align=center, above] {$r$} ($(device.east) + (1, -4.4)$);

% Server
\node [align=center] at ($(server.south) + (0, -3.4)$) {
    $s \gets {\sf Compare}({\sf csk}, \mathbf{c_x}, \mathbf{c_y})$ \\
    $r \gets {\sf Verify}(s)$
};


\end{tikzpicture}

\caption{Usage Model on Authentication}
\label{fig:model_auth}
\end{figure}



%-------------------

For example, let ${\sf FE} = ({\sf FE.Setup}, {\sf FE.KeyGen}, {\sf FE.Enc}, {\sf FE.Dec})$ be an fh-IPFE scheme we defined in Definition \ref{def:fh-IPFE}. Following \cite{cryptoeprint:2023/481}, we can instantiate a biometric authentication scheme using ${\sf FE}$ with the distance metric the Euclidean distance. Let the biometric templates $\mathbf{b}$ and $\mathbf{b'}$ be sampled from some distribution $\mathcal{B} \subseteq [m]^k$, and let the associated field of ${\sf FE}$ be $\mathbb{Z}_q$ where $q$ is a prime number larger than the maximum possible Euclidean distance $m^2 \cdot k$. The scheme is instantiated as follows.

\begin{itemize}

	\item ${\sf USetup}(1^\lambda)$: It calls ${\sf FE.Setup}(1^\lambda) \to {\sf msk}, {\sf pp}$ and outputs $({\sf esk}, {\sf psk}) \gets (({\sf msk}, {\sf pp}), ({\sf msk}, {\sf pp}))$ and ${\sf csk} \gets {\sf pp}$.

	\item ${\sf encodeEnroll}^{\mathcal{O}_{\mathcal{B}}}() \to \mathbf{x}$: For a template vector $\mathbf{b} = (b_1, b_2, \cdots, b_k)$ sampled from $\mathcal{O}_{\mathcal{B}}$, the function encodes it as $\mathbf{x} = (x_1, x_2, \cdots, x_{k+2}) = (b_1, b_2, \cdots, b_k, 1, \|\mathbf{b}\|^2)$. The auxiliary oracle is empty.

	\item ${\sf Enroll}({\sf esk}, \mathbf{x})$: It calls ${\sf FE.KeyGen}({\sf msk}, {\sf pp}, \mathbf{x}) \to f_\mathbf{x}$ and outputs $\mathbf{c_x} \gets f_\mathbf{x}$.

	\item ${\sf encodeProbe}^{\mathcal{O}_{\mathcal{B}}}()$: For a template vector $\mathbf{b'} = (b_1', b_2', \cdots, b_k')$ sampled from $\mathcal{O}_{\mathcal{B}}$, the function encodes it as $\mathbf{y} = (y_1, y_2, \cdots, y_{k+2}) = (-2b_1', -2b_2', \cdots, -2b_k', \|\mathbf{b'}\|^2, 1)$. The auxiliary oracle is empty.

	\item ${\sf Probe}({\sf psk}, \mathbf{y})$: It calls ${\sf FE.Enc}({\sf msk}, {\sf pp}, \mathbf{y}) \to \mathbf{c_y}$ and outputs $\mathbf{c_y}$.

	\item ${\sf Compare}({\sf csk}, \mathbf{c_x}, \mathbf{c_y)}$: It calls ${\sf FE.Dec}({\sf pp}, \mathbf{c_x}, \mathbf{c_y}) \to s$ and outputs the value $s$.

	\item ${\sf Verify}(s)$: If $\sqrt{s} < \tau$, a pre-defined threshold for comparing the closeness of two templates, then it outputs $r = 1$; otherwise, it outputs $r = 0$.

\end{itemize}

By the correctness of the functional encryption scheme $\sf FE$, we have
\[
	s = {\sf FE.Dec}({\sf pp}, \mathbf{c_x}, \mathbf{c_y}) = \langle \mathbf{x}, \mathbf{y} \rangle = \sum_{i=1}^k -2b_ib_i' + \|\mathbf{b}\|^2 + \|\mathbf{b'}\|^2 = \| \mathbf{b} - \mathbf{b'} \|^2.
\]

which is the square of the Euclidean distance between two templates $\mathbf{b}$ and $\mathbf{b}'$. Therefore, if two templates $\mathbf{b}$ and $\mathbf{b}'$ are close enough such that $\|\mathbf{b} - \mathbf{b'}\| < \tau$, the scheme results in $r = 1$, a successful authentication.


%-------------------

\section{Security Games}

From now on, we consider a family of biometric distributions $\mathbb{B}$. Removing a person $\mathcal{B}$ from $\mathbb{B}$ is written as $\mathbb{B} \setminus \mathcal{B}$.

\subsection{Forgery Game}
\label{sec:forgery_game}

In the forgery game, we model the ability of a malicious {\sf Scanner} who has access to the server's database of registered enrollments and tries to forge the user. The adversary $\mathcal{A}$ is given the enrollment message $\mathbf{c_x}$ and oracle $\mathcal{O}$ and tries to find a valid probe message $\mathbf{\tilde{z}}$. The whole game is defined in Figure \ref{fig:forgery_game}.

\begin{figure}[h]
\centerline{

\minipage{0.45\textwidth}
	\begin{tabular}{l c}
		${\sf Forg}_{\Pi, \mathbb{B}}(\mathcal{A}^\mathcal{O})$\\

			\hline

			$\mathcal{B} \getsdollar \mathbb{B}, \mathbb{B} \gets \mathbb{B} \setminus \mathcal{B}$ \\

			${\sf esk}, {\sf psk}, {\sf csk} \gets {\sf USetup}(1^\lambda)$ \\

			$\mathbf{x} \gets {\sf encodeEnroll}^{\mathcal{O}_{\mathcal{B}}}()$ \\

			$\mathbf{c_x} \gets {\sf Enroll}({\sf esk}, \mathbf{x})$ \\

			${\mathbf{\tilde{z}}} \gets \mathcal{A}^{ \mathcal{O} } ( \mathbf{c_x} )$ \\	

			$s \gets {\sf Compare}( {\sf csk}, \mathbf{c_x}, \mathbf{\tilde{z}} )$ \\

			\textbf{return} ${\sf Verify}(s)$
			
	\end{tabular}
	\caption{The Forgery Game}
	\label{fig:forgery_game}
\endminipage
\minipage{0.45\textwidth}
	\begin{tabular}{l c}
		${\sf Forg'}_{\Pi, \mathbb{B}}(\mathcal{A'}^{\mathcal{O'}})$\\

			\hline

			$\mathcal{B} \getsdollar \mathbb{B}, \mathbb{B} \gets \mathbb{B} \setminus \mathcal{B}$ \\

			${\sf esk}, {\sf psk}, {\sf csk} \gets {\sf USetup}(1^\lambda)$ \\

			$\mathbf{x} \gets {\sf encodeEnroll}^{\mathcal{O}_{\mathcal{B}}}()$ \\

			$\mathbf{c_x} \gets {\sf Enroll}({\sf esk}, \mathbf{x})$ \\

			${\mathbf{\tilde{z}}} \gets \mathcal{A'}^{\mathcal{O'}}()$ \\

			$s \gets {\sf Compare}( {\sf csk}, \mathbf{c_x}, \mathbf{\tilde{z}} )$ \\

			\textbf{return} ${\sf Verify}(s)$
			
		\end{tabular}
	\caption{The Plain Forgery Game}
	\label{fig:plain_forgery_game}
\endminipage
}
\end{figure}


The given oracle $\mathcal{O}$ consists of the following oracles:

\begin{itemize}
	\item $\mathcal{O}_{\mathcal{B}}$: It outputs a biometric sample $\mathbf{b} \getsdollar \mathcal{B}$.

	\item $\mathcal{O}_{\sf samp}(\cdot)$: On input an index $i$,
	\begin{itemize}
		\item If $i$ was not queried before, it first samples a biometric distribution $\mathcal{B}_i \in \mathbb{B}$ and then outputs a biometric sample $\mathbf{b} \getsdollar \mathcal{B}_i$.
		\item If $i$ has been queried before, it outputs a biometric sample $\mathbf{b} \getsdollar \mathcal{B}_i$.
	\end{itemize}

\end{itemize}

Depending on the threat model, we can also consider the following oracles:

\begin{itemize}
	\item $\mathcal{O}_{\sf Enroll}({\sf esk}, \cdot)$: On input $\mathbf{x'}$, it outputs the enrollment message ${\sf Enroll}({\sf esk}, \mathbf{x'})$.

	\item $\mathcal{O}_{\sf Probe}({\sf psk}, \cdot)$: On input $\mathbf{y'}$, it outputs the probe message ${\sf Probe}({\sf psk}, \mathbf{y'})$.

	\item $\mathcal{O}_{\sf Compare}^q({\sf csk}, \cdot, \cdot)$: If this oracle has been queried over $q$ times, it aborts. Otherwise, on input $\mathbf{c_x'}$ and $\mathbf{c_y'}$, it outputs the comparison result ${\sf Compare}( {\sf csk}, \mathbf{c_x'}, \mathbf{c_y'} )$.

\end{itemize}

Note that if the enrollment secret key ${\sf esk}$, probe secret key ${\sf psk}$, or the comparison secret key ${\sf csk}$ is an empty or public string in the scheme, then the corresponding oracles are naturally and implicitly given since the adversary can compute them herself. 

To consider potential false positives of biometrics match, we consider the plain forgery game in Figure \ref{fig:plain_forgery_game}, in which the adversary does not have any knowledge about the template. The given oracle $\mathcal{O'}$ consists of:
\begin{itemize}
	\item $\mathcal{O}_{\sf samp}(\cdot)$: On input an index $i$,
	\begin{itemize}
		\item If $i$ was not queried before, it first samples a biometric distribution $\mathcal{B}_i \in \mathbb{B}$ and then outputs a biometric sample $\mathbf{b} \getsdollar \mathcal{B}_i$.
		\item If $i$ has been queried before, it outputs a biometric sample $\mathbf{b} \getsdollar \mathcal{B}_i$.
	\end{itemize}
	
	\item $\mathcal{O}_{\sf auth}({\sf csk}, \mathbf{c_x}, \cdot)$: On input $\mathbf{z}$, it outputs ${\sf Verify}({\sf Compare}({\sf csk}, \mathbf{c_x}, \mathbf{z}))$.
\end{itemize}


We define the advantage of an adversary $\mathcal{A}$ in the forgery game of a scheme $\Pi$ associated with a family of distributions $\mathbb{B}$ as
\[
	\Adv^{\sf Forg}_{\Pi, \mathbb{B}, \mathcal{A}^\mathcal{O}} := \Pr[{\sf Forg}_{\Pi, \mathbb{B}}(\mathcal{A}^\mathcal{O}) \to 1] -
	\max_{\text{PPT } \mathcal{A'}} \Pr[{\sf Forg'}_{\Pi, \mathbb{B}}(\mathcal{A'}^\mathcal{O'}) \to 1].
\]

An authentication scheme $\Pi$ associated with a family of distributions $\mathbb{B}$ is called \emph{forgery secure} if for any PPT adversary $\mathcal{A}$,
\[
	\Adv^{\sf Forg}_{\Pi, \mathbb{B}, \mathcal{A^\mathcal{O}}} = \negl.
\]

%-------------------

\subsection{Simulation Game}
\label{sec:simulation_game}

In the simulation game, we model the ability of the {\sf Server} who tries to learn something more than the comparison result of the enrollment and probe messages. The adversary $\mathcal{A}$ is given an enrollment message and a list of $t$ probe messages, and it needs to guess whether these are real messages or simulation results of a simulator $\mathcal{S} = (\mathcal{S}_0, \mathcal{S}_{\sf Enroll}, \mathcal{S}_{\sf Probe})$ based on the $\sf Compare$ function. Intuitively, the simulator receives a list of $\sf Compare$ results of real enrollment and probe messages, and it returns some manual enrollment or probe messages that look similar to real ones. The whole game is defined in Figure \ref{fig:sim_game}.

\begin{figure}[h]
	\begin{center}

		\begin{subfigure}[t]{0.49\textwidth}
		\begin{tabular}{l c}
			${\sf SIM-Real}_{\Pi, \mathbb{B}}(\mathcal{A}^\mathcal{O})$\\

			\hline

			$\mathcal{B} \getsdollar \mathbb{B}, \mathbb{B} \gets \mathbb{B} \setminus \mathcal{B}$ \\

			${\sf esk}, {\sf psk}, {\sf csk} \gets {\sf USetup}(1^\lambda)$ \\

			$\mathbf{x} \gets {\sf encodeEnroll}^{\mathcal{O}_{\mathcal{B}}}()$ \\

			$\mathbf{c_x} \gets {\sf Enroll}({\sf esk}, \mathbf{x})$ \\

			$\{\mathbf{y}^{(i)}\}_{i=1}^t \gets \{{\sf encodeProbe}^{\mathcal{O}_{\mathcal{B}}}() \}_{i=1}^t $ \\
			
			$\{ \mathbf{c_y}^{(i)} \}_{i=1}^t \gets \left\{ {\sf Probe}( {\sf psk}, \mathbf{y}^{(i)} ) \right \}_{i=1}^t$ \\

			$b \gets \mathcal{A}^\mathcal{O} ({\sf csk}, \mathbf{c_x}, \{\mathbf{c_y}^{(i)}\}_{i=1}^t )$ \\

			\textbf{return} $b$
			
		\end{tabular}
		\end{subfigure}
		\begin{subfigure}[t]{0.49\textwidth}
		\begin{tabular}{l c}
			${\sf SIM-Ideal}_{\Pi, \mathbb{B}}(\mathcal{A}^{\tilde{\mathcal{O}}}, \mathcal{S} )$\\

			\hline

			$\mathcal{B} \getsdollar \mathbb{B}, \mathbb{B} \gets \mathbb{B} \setminus \mathcal{B}$ \\

			${\sf esk}, {\sf psk}, {\sf csk}, \gets {\sf USetup}(1^\lambda)$ \\

			$ \mathbf{x} \gets {\sf encodeEnroll}^{\mathcal{O}_{\mathcal{B}}}()$ \\

			$\mathbf{c_x} \gets {\sf Enroll}({\sf esk}, \mathbf{x})$ \\

			$\{\mathbf{y}^{(i)}\}_{i=1}^t \gets \{{\sf encodeProbe}^{\mathcal{O}_{\mathcal{B}}}()\}_{i=1}^t $ \\
			
			$\{ \mathbf{c_y}^{(i)} \}_{i=1}^t \gets \left\{ {\sf Probe}( {\sf psk}, \mathbf{y}^{(i)} ) \right \}_{i=1}^t$ \\

			${\sf C} \gets \{ {\sf Compare}({\sf csk}, \mathbf{c_x}, \mathbf{c_y}^{(i)}) \}_{i=1}^t$ \\

			$\mathbf{\tilde c_x}, \{\mathbf{\tilde c_y}^{(i)}\}_{i=1}^t \gets \mathcal{S}_0( {\sf csk}, {\sf C} )$ \\

			$b \gets \mathcal{A}^\mathcal{\tilde O} ({\sf csk}, \mathbf{\tilde c_x}, \{ \mathbf{\tilde c_y}^{(i)} \}_{i=1}^t )$ \\

			\textbf{return} $b$
			
		\end{tabular}
		\end{subfigure}
	\end{center}
	\caption{The Simulation Game}
	\label{fig:sim_game}
\end{figure}

The given oracle $\mathcal{O}$ and $\mathcal{\tilde O}$ consist of the oracle $\mathcal{O}_{\sf samp}$:

\begin{itemize}
	\item $\mathcal{O}_{\sf samp}(\cdot)$: On input an index $i$,
	\begin{itemize}
		\item If $i$ was not queried before, it first samples a biometric distribution $\mathcal{B}_i \in \mathbb{B}$ and then outputs a biometric sample $\mathbf{b} \getsdollar \mathcal{B}_i$.
		\item If $i$ has been queried before, it outputs a biometric sample $\mathbf{b} \getsdollar \mathcal{B}_i$.
	\end{itemize}

\end{itemize}

Depending on the threat model, we can also consider the following oracles for $\mathcal{O}$:

\begin{itemize}
	\item $\mathcal{O}_{\sf Enroll}({\sf esk}, \cdot)$: On input $\mathbf{x'}$, it outputs the enrollment message ${\sf Enroll}({\sf esk}, \mathbf{x'})$.

	\item $\mathcal{O}_{\sf Probe}({\sf psk}, \cdot)$: On input $\mathbf{y'}$, it outputs the probe message ${\sf Probe}({\sf psk}, \mathbf{y'})$.
\end{itemize}

Their corresponding oracles for $\mathcal{\tilde O}$ are stateful and memorize all the previous queries. It includes two simulators: $\mathcal{S}_{\sf Enroll}$ and $\mathcal{S}_{\sf Probe}$.

\begin{itemize}

	\item $\tilde{\mathcal{O}}_{\sf Enroll} ({\sf csk} ,{\sf esk}, \cdot)$: On input $\mathbf{x}^{(j)}$, it updates the collection of comparison results ${\sf C}$ by adding the results with all previous probe messages $\mathbf{c_y}^{(i)}$. Then it calls the simulator $\mathcal{S}_{\sf Enroll}({\sf csk}, {\sf C})$ and returns whatever the simulator returns.

	\item $\tilde{\mathcal{O}}_{\sf Probe} ({\sf csk}, {\sf psk}, \cdot)$: On input $\mathbf{y}^{(j)}$, it updates the collection of comparison results ${\sf C}$ by adding the results with all previous enrollment messages $\mathbf{c_x}^{(i)}$. Then it calls the simulator $\mathcal{S}_{\sf Probe}({\sf csk}, {\sf C})$ and returns whatever the simulator returns.

\end{itemize}

The details of these oracles are given in Figure \ref{fig:sim_game_oracle}.

\begin{figure}[h]
	\begin{center}

		\begin{subfigure}[t]{0.49\textwidth}
		\begin{tabular}{l c}
			${\mathcal{ \tilde O}}_{\sf Enroll}({\sf csk}, {\sf esk}, \mathbf{x}^{(j)})$\\

			\hline

			$\mathbf{c_x}^{(j)} \gets {\sf Enroll}({\sf esk}, \mathbf{x}^{(j)})$ \\

			${\sf C} \gets {\sf C} \cup \{ {\sf Compare}({\sf csk}, \mathbf{c_x}^{(j)}, \mathbf{c_y}^{(i)}) \}_{i}$ \\

			$\mathbf{c_x}' \gets \mathcal{S}_{\sf Enroll}({\sf csk}, {\sf C} )$ \\
			
			\textbf{return} $\mathbf{c_x}'$
			
		\end{tabular}
		\end{subfigure}
		\begin{subfigure}[t]{0.49\textwidth}
		\begin{tabular}{l c}
			$\tilde{\mathcal{O}}_{\sf Probe}({\sf csk}, {\sf psk}, \mathbf{y}^{(j)})$\\

			\hline

			$\mathbf{c_y}^{(j)} \gets {\sf Probe}( {\sf psk}, \mathbf{y}^{(j)} ) $ \\

			${\sf C} \gets {\sf C} \cup \{ {\sf Compare}({\sf csk}, \mathbf{c_x}^{(i)}, \mathbf{c_y}^{(j)}) \}_{i}$ \\

			$\mathbf{c_y}' \gets \mathcal{S}_{\sf Probe}({\sf csk}, {\sf C} )$ \\

			\textbf{return} $\mathbf{c_y}'$
			
		\end{tabular}
		\end{subfigure}
	\end{center}
	\caption{Choice of the Oracle $\tilde{\mathcal{O}}$}
	\label{fig:sim_game_oracle}
\end{figure}


We define the advantage of an adversary $\mathcal{A}$ and a simulator $\mathcal{S} = (\mathcal{S}_0, \mathcal{S}_{\sf Enroll}, \mathcal{S}_{\sf Probe})$ in the simulation game of a scheme $\Pi$ associated with a family of distributions $\mathbb{B}$ as
\[
	\Adv^{\sf SIM}_{\Pi, \mathbb{B}, \mathcal{A}^{\mathcal{O}, \tilde{\mathcal{O}}}, \mathcal{S}} := |\Pr[{\sf SIM-Real}_{\Pi, \mathbb{B}}(\mathcal{A}^\mathcal{O}) \to 1] - \Pr[{\sf SIM-Ideal}_{\Pi, \mathbb{B}}(\mathcal{A}^{\tilde{\mathcal{O}}}, \mathcal{S}) \to 1]|.
\]

An authentication scheme $\Pi$ associated with a family of distributions $\mathbb{B}$ is called \emph{simulation secure} if for any PPT adversary $\mathcal{A}$, there exists a PPT simulator $\mathcal{S} = (\mathcal{S}_0, \mathcal{S}_{\sf Enroll}, \mathcal{S}_{\sf Probe})$ such that
\[
	\Adv^{\sf SIM}_{\Pi, \mathbb{B}, \mathcal{A}^{\mathcal{O}, \tilde{\mathcal{O}}}, \mathcal{S}} = \negl.
\]

%-------------------

\subsection{Identification Game}
\label{sec:id_game}

In the identification game, we model the ability of the {\sf Server} who tries to identify the user. The adversary $\mathcal{A}$ is given an enrollment message $\mathbf{c_x}$, a list of $t$ probe messages $\{ \mathbf{c_y}^{(i)} \}_{i=1}^t$, and oracles to two distinct distributions $\mathcal{B}^{(0)}, \mathcal{B}^{(1)}$. It tries to guess from which these messages are generated. The whole game is defined in Figure \ref{fig:id_game}.

\begin{figure}[h]
	\begin{center}
	\begin{tabular}{l c}
		${\sf Id}_{\Pi, \mathbb{B}}(\mathcal{A}^\mathcal{O})$\\

			\hline

			$b \getsdollar \{0, 1\}$ \\

			$\mathcal{B}^{(0)}, \mathcal{B}^{(1)} \getsdollar \mathbb{B}, \mathbb{B} \gets \mathbb{B} \setminus \{ \mathcal{B}^{(0)}, \mathcal{B}^{(1)} \}$ \\

			${\sf esk}, {\sf psk}, {\sf csk}, \gets {\sf USetup}(1^\lambda)$ \\

			$\mathbf{x} \gets {\sf encodeEnroll}^{\mathcal{O}_{\mathcal{B}^{(b)}}}()$ \\

			$\mathbf{c_x} \gets {\sf Enroll}({\sf esk}, \mathbf{x})$ \\

			$\{\mathbf{y}^{(i)}\}_{i=1}^t \gets \{ {\sf encodeProbe}^{\mathcal{O}_{\mathcal{B}^{(b)}}}() \}_{i=1}^t $ \\
			
			$\{ \mathbf{c_y}^{(i)} \}_{i=1}^t \gets \left\{ {\sf Probe}( {\sf psk}, \mathbf{y}^{(i)} ) \right \}_{i=1}^t$ \\

			$\tilde{b} \gets \mathcal{A}^\mathcal{O} ( {\sf csk}, \mathbf{c_x}, \{ \mathbf{c_y}^{(i)} \}_{i=1}^t )$ \\

			\textbf{return} $1_{\tilde{b} = b}$
			
	\end{tabular}
	\end{center}
	\caption{The Identification Game}
	\label{fig:id_game}
\end{figure}


The given oracle $\mathcal{O}$ consists of the following oracles:

\begin{itemize}
	\item $\mathcal{O}_{\mathcal{B}^{(0)}}$: It outputs a biometric sample $\mathbf{b} \getsdollar \mathcal{B}^{(0)}$.

	\item $\mathcal{O}_{\mathcal{B}^{(1)}}$: It outputs a biometric sample $\mathbf{b} \getsdollar \mathcal{B}^{(1)}$.

	\item $\mathcal{O}_{\sf samp}(\cdot)$: On input an index $i$,
	\begin{itemize}
		\item If $i$ was not queried before, it first samples a biometric distribution $\mathcal{B}_i \in \mathbb{B}$ and then outputs a biometric sample $\mathbf{b} \getsdollar \mathcal{B}_i$.
		\item If $i$ has been queried before, it outputs a biometric sample $\mathbf{b} \getsdollar \mathcal{B}_i$.
	\end{itemize}
	
\end{itemize}

Depending on the threat model, we can also consider the oracle:

\begin{itemize}
	\item $\mathcal{O}_{\sf Enroll}({\sf esk}, \cdot)$: On input $\mathbf{x'}$, it outputs the enrollment message ${\sf Enroll}({\sf esk}, \mathbf{x'})$.
\end{itemize}

We define the advantage of an adversary $\mathcal{A}$ in the identification game of a scheme $\Pi$ associated with a family of distributions $\mathbb{B}$ as
\[
	\Adv^{\sf Id}_{\Pi, \mathbb{B}, \mathcal{A}^\mathcal{O}} := |\Pr[{\sf Id}_{\Pi}(\mathcal{A}^\mathcal{O}) \to 1] - \frac{1}{2}|.
\]

An authentication scheme $\Pi$ associated with a family of distributions $\mathbb{B}$ is called \emph{identification secure} if for any PPT adversary $\mathcal{A}$,
\[
	\Adv^{\sf Id}_{\Pi, \mathbb{B}, \mathcal{A}^\mathcal{O}} = \negl.
\]

%-------------------

\subsection{Reusability Game}
\label{sec:reuse_game}

In the reusability game, we model the ability of a malicious application who tries to forge the user in another application. The adversay $\mathcal{A}$ is given the enrollment message $\mathbf{c_x}$ and oracle $\mathcal{O}$ and tries to find a valid probe message $\mathbf{\tilde{z}}$. The whole game is defined in Figure \ref{fig:reuse_game}.

\begin{figure}[h]
\centerline{

\minipage{0.45\textwidth}
	\begin{tabular}{l c}
		${\sf Reuse}_{\Pi, \mathbb{B}}(\mathcal{A}^\mathcal{O})$\\

			\hline

			$\mathcal{B} \getsdollar \mathbb{B}, \mathbb{B} \gets \mathbb{B} \setminus \mathcal{B}$ \\

			${\sf esk}, {\sf psk}, {\sf csk} \gets {\sf USetup}(1^\lambda)$ \\

			$\mathbf{x} \gets {\sf encodeEnroll}^{\mathcal{O}_{\mathcal{B}}}()$ \\

			$\mathbf{c_x} \gets {\sf Enroll}({\sf esk}, \mathbf{x})$ \\

			${\mathbf{\tilde{z}}} \gets \mathcal{A}^{\mathcal{O}} ( \mathbf{c_x} )$ \\

			$s \gets {\sf Compare}( {\sf csk}, \mathbf{c_x}, \mathbf{\tilde{z}} )$ \\

			\textbf{return} ${\sf Verify}(s)$
			
	\end{tabular}
	\caption{The Reusability Game}
	\label{fig:reuse_game}
\endminipage
\minipage{0.48\textwidth}
	\begin{tabular}{l c}
		${\sf Reuse'}_{\Pi, \mathbb{B}}(\mathcal{A'}^{\mathcal{O'}})$\\

			\hline

			$\mathcal{B} \getsdollar \mathbb{B}, \mathbb{B} \gets \mathbb{B} \setminus \mathcal{B}$ \\

			${\sf esk}, {\sf psk}, {\sf csk} \gets {\sf USetup}(1^\lambda)$ \\

			$\mathbf{x} \gets {\sf encodeEnroll}^{\mathcal{O}_{\mathcal{B}}}()$ \\

			$\mathbf{c_x} \gets {\sf Enroll}({\sf esk}, \mathbf{x})$ \\

			${\mathbf{\tilde{z}}} \gets \mathcal{A'}^{\mathcal{O'}}()$ \\

			$s \gets {\sf Compare}( {\sf csk}, \mathbf{c_x}, \mathbf{\tilde{z}} )$ \\

			\textbf{return} ${\sf Verify}(s)$
			
		\end{tabular}
	\caption{The Plain Reusability Game}
	\label{fig:plain_reuse_game}
\endminipage
}
\end{figure}

The given oracle $\mathcal{O}$ consist of the following oracles:

\begin{itemize}
	\item $\mathcal{O}_{\sf samp}(\cdot)$: On input an index $i$,
	\begin{itemize}
		\item If $i$ was not queried before, it first samples a biometric distribution $\mathcal{B}_i \in \mathbb{B}$ and then outputs a biometric sample $\mathbf{b} \getsdollar \mathcal{B}_i$.
		\item If $i$ has been queried before, it outputs a biometric sample $\mathbf{b} \getsdollar \mathcal{B}_i$.
	\end{itemize}
	
\item $\mathcal{O}_{\sf Enroll}^{ {\sf Re}} (\cdot)$: On input ${\sf esk'}$, it first samples $\mathbf{x'} \getsdollar {\sf encodeEnroll}^{\mathcal{O}_{\mathcal{B}}}()$ and outputs ${\sf Enroll}({\sf esk'}, \mathbf{x'})$.
	
	\item $\mathcal{O}_{\sf Probe}^{ {\sf Re}} (\cdot)$: On input ${\sf psk'}$, it first samples $\mathbf{y'} \getsdollar {\sf encodeProbe}^{\mathcal{O}_{\mathcal{B}}}()$ and outputs ${\sf Probe}({\sf psk'}, \mathbf{y'})$.

\end{itemize}

Depending on the threat model, we can also consider the following oracles:

\begin{itemize}
	\item $\mathcal{O}_{\sf Enroll}({\sf esk}, \cdot)$: On input $\mathbf{x'}$, it outputs the enrollment message ${\sf Enroll}({\sf esk}, \mathbf{x'})$.

	\item $\mathcal{O}_{\sf Probe}({\sf psk}, \cdot)$: On input $\mathbf{y'}$, it outputs the probe message ${\sf Probe}({\sf psk}, \mathbf{y'})$.

	\item $\mathcal{O}_{\sf Compare}^q({\sf csk}, \cdot, \cdot)$: If this oracle has been queried over $q$ times, it aborts. Otherwise, on input $\mathbf{c_x'}$ and $\mathbf{c_y'}$, it outputs the comparison result ${\sf Compare}( {\sf csk}, \mathbf{c_x'}, \mathbf{c_y'} )$.

\end{itemize}

Note that the resuability game is defined in a similar way as the forgery game in Section \ref{sec:forgery_game}. The difference is their use cases and oracles. In the forgery game, the targeted adversary is an eavesdropper who has access to the server's database. If it can also access the user's device, it may be able to execute $\mathcal{O}_{\sf Enroll}$ and $\mathcal{O}_{\sf Probe}$. In contrast, the reusability game targets a malicious application, who can ask the user to enroll and probe with crafted secret keys. We assume the malicious application cannot access other functions in other applications.

To consider potential false positives of biometrics match, we consider the plain reusability game in Figure \ref{fig:plain_reuse_game}, in which the adversary does not have any knowledge about the template. The given oracle $\mathcal{O'}$ consists of:

\begin{itemize}

	\item $\mathcal{O}_{\sf samp}(\cdot)$: On input an index $i$,
	\begin{itemize}
		\item If $i$ was not queried before, it first samples a biometric distribution $\mathcal{B}_i \in \mathbb{B}$ and then outputs a biometric sample $\mathbf{b} \getsdollar \mathcal{B}_i$.
		\item If $i$ has been queried before, it outputs a biometric sample $\mathbf{b} \getsdollar \mathcal{B}_i$.
	\end{itemize}
	
	\item $\mathcal{O}_{\sf auth}({\sf csk}, \mathbf{c_x}, \cdot)$: On input $\mathbf{z}$, it outputs ${\sf Verify}({\sf Compare}({\sf csk}, \mathbf{c_x}, \mathbf{z}))$.

\end{itemize}


We define the advantage of an adversary $\mathcal{A}$ in the reusability game of a scheme $\Pi$ associated with a family of distributions $\mathbb{B}$ as
\[
	\Adv^{\sf Reuse}_{\Pi, \mathbb{B}, \mathcal{A}^\mathcal{O}} := \Pr[{\sf Reuse}_{\Pi, \mathbb{B}}(\mathcal{A}^\mathcal{O}) \to 1] -
	\max_{\text{PPT } \mathcal{A}'} \Pr[{\sf Reuse'}_{\Pi, \mathbb{B}}(\mathcal{A'}^\mathcal{O'}) \to 1].
\]

An authentication scheme $\Pi$ associated with a family of distributions $\mathbb{B}$ is called \emph{reusability secure} if for any PPT adversary $\mathcal{A}$,
\[
	\Adv^{\sf Reuse}_{\Pi, \mathbb{B}, \mathcal{A^\mathcal{O}}} = \negl.
\]

%-------------------

%\section{Analysis of Games}

\subsection{Id Games $\leq$ Forgery Games}

We provide a reduction to show that, given the following assumption, if an adversary $\mathcal{A}_{\sf Forg}$ can break the forgery game with oracles $\mathcal{O}_{\sf samp}$, $\mathcal{O}_{\sf Enroll}$ and $\mathcal{O}_{\sf Compare}$, then there exists an adversary $\mathcal{A}_{\sf Id}$ who can break the identification game if it is given the oracle $\mathcal{O}_{\sf samp}$ and $\mathcal{O}_{\sf Enroll}$.

\begin{assumption}
	\label{assump:id_lt_forg}
	Let ${\sf esk}, {\sf psk}, {\sf csk} \gets {\sf USetup}(1^\lambda)$, $\mathbf{x}^{(0)}, \mathbf{x}^{(1)} \gets {\sf encodeEnroll}^{\mathcal{O}_{\mathcal{B}}}()$, $\mathbf{c_x}^{(0)} \gets {\sf Enroll}({\sf esk}, \mathbf{x}^{(0)})$, $ \mathbf{c_x}^{(1)} \gets {\sf Enroll}({\sf esk}, \mathbf{x}^{(1)})$. Then there exists a constant $\tau$ such that for all $\mathbf{z}$ where
	\[
		{\sf Verify}( {\sf Compare}({\sf csk}, \mathbf{c_x}^{(0)}, \mathbf{z}) ) = 1
	\]
	we have
	\[
		\Pr[{\sf Verify}( {\sf Compare}({\sf csk}, \mathbf{c_x}^{(1)}, \mathbf{z}) ) = 1 ] \geq \tau.
	\]	

\end{assumption}

\begin{theorem}

	Given Assumption \ref{assump:id_lt_forg} with a constant $\tau$. Let $\mathcal{A}_{\sf Forg}^{\mathcal{O}_{\sf Forg}}$ be an adversary in the Forgery game with oracles $\mathcal{O}_{\sf Forg} = (\mathcal{O'}_{\sf samp}, \mathcal{O'}_{\sf Enroll}, \mathcal{O'}_{\sf Compare})$. Then there exists an adversary $\mathcal{A}_{\sf Id}$ in the identification game with oracles $\mathcal{O}_{\sf Id} = (\mathcal{O}_{\sf samp}, \mathcal{O}_{\sf Enroll})$ such that
	\[
		\Adv_{\Pi, \mathbb{B}, \mathcal{A}_{\sf Id}^{\mathcal{O}_{\sf Id}}} \geq \frac{1}{4} [ \Adv_{\Pi, \mathbb{B}, \mathcal{A}_{\sf Forg}^{\mathcal{O}_{\sf Forg}}}  - (1-\tau) ].
	\]

\end{theorem}

\begin{proof}

	Consider the adversary $\mathcal{A}_{\sf Id}$ in Figure \ref{fig:id_game_ad} with oracles $\mathcal{O}_{\sf Id} = (\mathcal{O}_{\sf samp}, \mathcal{O}_{\sf Enroll})$.

\begin{figure}[htp]
	\begin{center}
	\begin{tabular}{l c}
		$\mathcal{A}_{\sf Id}^{\mathcal{O}_{\sf Id}}({\sf csk}, \mathbf{c_x}^{(b)}, \mathcal{B}^{(0)}, \mathcal{B}^{(1)})$\\

			\hline

			$\mathbf{x'} \gets {\sf encodeEnroll}^{\mathcal{O}_{\mathcal{B}^{(0)}}}()$ \\

			\qquad When $\mathcal{O}_{\mathcal{B}^{(0)}}$ is called, sample $\mathbf{b} \gets \mathcal{B}^{(0)}$ and return $\mathbf{b}$. \\

			$\mathbf{c_x'} \gets \mathcal{O}_{\sf Enroll}(\mathbf{x'})$ \\

			$\mathbf{\tilde{z}} \gets \mathcal{A}_{\sf Forg}^{\mathcal{O}_{\sf Forg}} ( \mathbf{c_x'} )$ \\
			\qquad When $\mathcal{O'}_{\sf samp}(i)$ is called, return $\mathcal{O}_{\sf samp}(i)$ \\
			\qquad When $\mathcal{O'}_{\sf Enroll}(\mathbf{x})$ is called, return $\mathcal{O}_{\sf Enroll}(\mathbf{x})$. \\
			\qquad When $\mathcal{O'}_{\sf Compare}(\mathbf{c_x}, \mathbf{c_y})$ is called, return ${\sf Compare}({\sf csk}, \mathbf{c_x}, \mathbf{c_y})$. \\

			$s \gets {\sf Compare}({\sf csk}, \mathbf{c_x}^{(b)}, \mathbf{\tilde{z}})$ \\

			If ${\sf Verify}(s) = 1$ \\
				\qquad \textbf{return} $0$ \\

			Else \\
				\qquad \textbf{return} a random bit $b' \getsdollar \{0, 1\}$
			
	\end{tabular}
	\end{center}
	\caption{The Identification Game Adversary}
	\label{fig:id_game_ad}
\end{figure}

Note that $\mathbf{c_x}'$ and $\mathcal{O'}_{\sf Enroll}, \mathcal{O'}_{\sf Compare}$ are simulated perfectly as the real forgery game with disbribution $\mathcal{B}^{(0)}$, so
\[
	\Pr[{\sf Verify}( {\sf Compare}({\sf csk}, \mathbf{c_x}', \mathbf{\tilde{z}}) ) \to 1] = \Pr[{\sf Forg}_{\Pi, \mathbb{B}}(\mathcal{A}_{\sf Forg}^{\mathcal{O}_{\sf Forg} }) \to 1 ].
\]
Therefore,
\[
	\Pr[{\sf Verify}( {\sf Compare}({\sf csk}, \mathbf{c_x}^{(0)}, \mathbf{\tilde{z}}) ) \to 1] \geq \Pr[{\sf Forg}_{\Pi, \mathbb{B}}(\mathcal{A}_{\sf Forg}^{\mathcal{O}_{\sf Forg}}) \to 1 ] \cdot \tau.
\]

Let the challenge bit be $b$. If $b = 0$, the probability that $\mathcal{A}_{\sf Id}^{\mathcal{O}_{\sf Enroll}}$ returns $b$ is
	\[
	\begin{aligned}
		&\Pr[{\sf Verify}({\sf Compare}({\sf csk}, \mathbf{c_x}^{(0)}, \mathbf{\tilde{z}})) \to 1 ] + (1 - \Pr[{\sf Verify}({\sf Compare}({\sf csk}, \mathbf{c_x}^{(0)}, \mathbf{\tilde{z}})) \to 1 ]) \cdot \frac{1}{2} \\
		= & \frac{1}{2} + \frac{1}{2}\Pr[{\sf Verify}({\sf Compare}({\sf csk}, \mathbf{c_x}^{(0)}, \mathbf{\tilde{z}})) \to 1 ] \geq \frac{1}{2} + \frac{\tau}{2} \Pr[{\sf Forg}_{\Pi, \mathbb{B}}(\mathcal{A}_{\sf Forg}^{\mathcal{O}_{\sf Forg} }) \to 1 ].
	\end{aligned}
	\]

If $b = 1$, the probability that $\mathcal{A}_{\sf Id}^{\mathcal{O}_{\sf Enroll}}$ returns $b$ is
	\[
	\begin{aligned}
		&\Pr[{\sf Verify}({\sf Compare}({\sf csk}, \mathbf{c_x}^{(1)}, \mathbf{\tilde{z}})) \to 0 ] \cdot \frac{1}{2} \\
		= & (1 - \Pr[{\sf Verify}({\sf Compare}({\sf csk}, \mathbf{c_x}^{(1)}, \mathbf{\tilde{z}})) \to 1 ] ) \cdot \frac{1}{2} \\
		\end{aligned}
	\]
	Note that $\mathcal{A}_{\sf Forg}$ is not given any knowledge about $\mathcal{B}^{(1)}$. The distribution of $\mathbf{c_x'}$ is identical to $\mathbf{c_x'} \gets \mathcal{O'}_{\sf Enroll}( {\sf encodeEnroll}^{\mathcal{O}_{\mathcal{B'}}}() ) $ where $\mathcal{B}' \getsdollar \mathbb{B} $, which it can sample itself in the plain forgery game with oracles $\mathcal{O}_{\sf samp}$ and $\mathcal{O}_{\sf Enroll}$. Therefore,
	\[
		\Pr[{\sf Verify}({\sf Compare}({\sf csk}, \mathbf{c_x}^{(1)}, \mathbf{\tilde{z}})) \to 1 ] \leq \max_{\text{PPT }\mathcal{A'}} \Pr[ {\sf Forg'}_{\Pi, \mathbb{B}}(\mathcal{A'}_{\sf Forg}^{\mathcal{O}_{\sf Forg}}) \to 1 ].
	\]

In summary,
\begin{equation*}
	\begin{aligned}
		\Pr[{\sf Id}_{\Pi}(\mathcal{A}_{\sf Id}^{\mathcal{O}_{\sf Id}}) \to 1] 
		&= \frac{1}{2} \Pr[\mathcal{A}_{\sf Id}^{\mathcal{O}_{\sf Id}} \to 0 \mid b = 0] + \frac{1}{2} \Pr[\mathcal{A}_{\sf Id}^{\mathcal{O}_{\sf Id}} \to 1 \mid b = 1] \\
		&\geq \frac{1}{4} + \frac{\tau}{4} \Pr[{\sf Forg}_{\Pi, \mathbb{B}}(\mathcal{A}_{\sf Forg}^{\mathcal{O}_{\sf Forg} }) \to 1 ] + \frac{1}{4} - \frac{1}{4}\max_{\text{PPT } \mathcal{A}'} \Pr[{\sf Forg'}_{\Pi, \mathbb{B}}(\mathcal{A'}_{\sf Forg}^{\mathcal{O}_{\sf Forg} }) \to 1 ] \\
		&= \frac{1}{2} + \frac{1}{4}\Adv_{\Pi, \mathbb{B}, \mathcal{A}_{\sf Forg}^{\mathcal{O}_{\sf Forg}}}  - \frac{1-\tau}{4}\Pr[{\sf Forg}_{\Pi, \mathbb{B}}(\mathcal{A}_{\sf Forg}^{\mathcal{O}_{\sf Forg} }) \to 1 ] \\
		&\geq \frac{1}{2} + \frac{1}{4}\Adv_{\Pi, \mathbb{B}, \mathcal{A}_{\sf Forg}^{\mathcal{O}_{\sf Forg}}}  - \frac{1-\tau}{4}
	\end{aligned}
\end{equation*}

\end{proof}





%-------------------

% Section for Attacks and Reduction
%1. If one can know whether x_1 is 0 or 1, it can break the fh-IPFE based scheme in the simuation game.
	%a. Show that distributions of <x,y> and <x,y> | x_1 = 0 are very close
	%b. Show that the simulator must follow the inner product rule

%2. If one can break the forgery game without E_probe, it can break the identification game with E_enroll and E_compare
	%a. First E_enroll(b^0)
	%b. Let the adversary for the forgery game finds a c_y
	%c. Check whether c_y is close to c_x^b or not
	%d. Need to specify that Compare(c_x^0, c_y^1) and Compare(c_x^1, c_y^0) are high

%3 For uniform distributed biometrics, if one can break the privacy game, one can break the simulation game.
	%a. Let the adversary for the privacy game on c_x return x1
	%b. x2 = x1 + r, r is uniformly generated
	%c. O_enroll(x2) returns cx2
	%d. Let the adversary for the privacy game on cx2 return x2'
	%e. Check whether x2' = x2
	%f. Show that in the real setting, x2' = x2 with high probability
	%g. Show that in the ideal setting, there are too many possibilities that cx2 can match the Compare results, so not likely a simulator who only knows the Compare results can return exactly x2'.

%-------------------
%% Backup File

% \section*{Backup}
\subsection*{Privacy Game}

In the privacy game, we model the ability of the server who tries to learn the biometric template of a user. The adversary $\mathcal{A}$ is given the enrollment message $\mathbf{c_x}$, $t$ probe messages $\{ \mathbf{c_y}_i\}_{i=1}^t$ for some integer $t$, and oracle $\mathcal{O}$ and tries to find the template $\mathbf{b}$. The whole game is defined in Figure \ref{fig:privacy_game}.


\begin{figure}[h]
	\begin{center}
		\begin{tabular}{l c}
			${\sf Priv}_{\Pi}(\mathcal{A})$\\

			\hline

			${\sf csk}, {\sf pp_s} \gets {\sf SSetup(1^\lambda)}$ \\

			${\sf esk}, {\sf psk}, {\sf pp_u} \gets {\sf USetup}(1^\lambda, {\sf pp_s})$ \\

			${\sf pp}:= ({\sf pp_s}, {\sf pp_u})$ \\

			$ \mathbf{b} \getsdollar \mathcal{B}, \mathbf{x} \gets {\sf encodeEnroll}^{\mathcal{O}_{\sf aux}}(\mathbf{b})$ \\

			$\mathbf{c_x} \gets {\sf Enroll}({\sf esk}, {\sf pp}, \mathbf{x})$ \\

			$\{ \mathbf{b}_i \}_{i=1}^t \getsdollar \mathcal{B}^t$ \\ 

			$\{ \mathbf{c_y}_i \}_{i=1}^t \gets \left\{ {\sf Probe}( {\sf psk}, {\sf pp}, {\sf encodeProbe}^{\mathcal{O}_{\sf aux}}(\mathbf{b}_i) ) \right \}_{i=1}^t$ \\

			${\mathbf{\tilde{b}}} \gets \mathcal{A}^{\mathcal{O}} ( \sf{pp}, \mathbf{c_x}, \{ \mathbf{c_y}_i \}_{i=1}^t )$ \\

			\textbf{return} $1_{\mathbf{\tilde{b}} = \mathbf{b}}$
			
		\end{tabular}
	\end{center}
	\caption{The Privacy Game}
	\label{fig:privacy_game}
\end{figure}


The given oracle $\mathcal{O}$ can be any or more of the following three oracles:


\begin{itemize}

	\item $\mathcal{O}_{\sf Enroll}({\sf esk}, {\sf pp}, \cdot)$: On input $\mathbf{x}$, it outputs the enrollment message ${\sf Enroll}({\sf esk}, {\sf pp}, \mathbf{x})$.

	\item $\mathcal{O}_{\sf Probe}({\sf psk}, {\sf pp}, \cdot)$: On input $\mathbf{y}$, it outputs the probe message ${\sf Probe}({\sf psk}, {\sf pp}, \mathbf{y})$.

	\item $\mathcal{O}_{\sf Compare}({\sf csk}, {\sf pp}, \cdot, \cdot)$: On input $\mathbf{c_x}$ and $\mathbf{c_y}$, it outputs the comparison result ${\sf Compare}( {\sf csk}, {\sf pp}, \mathbf{c_x}, \mathbf{c_y} )$.

\end{itemize}

Note that if the enrollment secret key ${\sf esk}$, the probe secret key ${\sf psk}$, or the comparison secret key ${\sf csk}$ is an empty string in the scheme, then the corresponding oracles are naturally and implicitly given since the adversary can compute them herself.

We define the success probability of an adversary $\mathcal{A}$ in the privacy game of the scheme $\Pi$ as
\[
	\Adv^{\sf Priv}_{\Pi, \mathcal{A}^\mathcal{O}} := \Pr[{\sf Priv}_{\Pi}(\mathcal{A}^\mathcal{O}) \to 1].
\]

The authentication scheme $\Pi$ is called \emph{privacy} secure if for any PPT adversary $\mathcal{A}$,
\[
	\Adv^{\sf Priv}_{\Pi, \mathcal{A}^\mathcal{O}} = \negl.
\]

%-------------------

\subsection*{IND Game}
For empty comparison key ${\sf csk}$, to avoid trivial attacks, we need to restrict the adversary's queries.

\begin{definition}[Admissible Adversary]
\label{def:admissible}
	For an adversary $\mathcal{A}^\mathcal{O}$ in an IND game of a scheme $\Pi$ with empty comparison secret key, let $(\mathbf{x}_0^1, \mathbf{x}_1^1), (\mathbf{x}_0^2, \mathbf{x}_1^2), \cdots, (\mathbf{x}_0^{Q_{\sf Enroll}}, \mathbf{x}_1^{Q_{\sf Enroll}})$ be the queries to $\mathcal{O}_{\sf Enroll}^{\sf IND}$, and let $(\mathbf{y}_0^1, \mathbf{y}_1^1), (\mathbf{y}_0^2, \mathbf{y}_1^2),$ $\cdots, (\mathbf{y}_0^{Q_{\sf Probe}}, \mathbf{y}_1^{Q_{\sf Probe}})$ be the queries to $\mathcal{O}_{\sf Probe}^{\sf IND}$. We say the adversary $\mathcal{A}$ is \emph{admissible} if
	\[
		\begin{aligned}
			& {\sf Compare}( {\sf pp}, {\sf Enroll}({\sf esk}, {\sf pp}, \mathbf{x}_0^i), {\sf Probe}({\sf psk}, {\sf pp}, \mathbf{y}_0^j) ) \\
			= \; & {\sf Compare}( {\sf pp}, {\sf Enroll}({\sf esk}, {\sf pp}, \mathbf{x}_1^i), {\sf Probe}({\sf psk}, {\sf pp}, \mathbf{y}_1^j) )
		\end{aligned}
		\quad \forall i \in [Q_{\sf Enroll}], \forall j \in [Q_{\sf Probe}]
	\]
	

\end{definition}

This prevents the adversary's win from simply querying unmatched enrollment and probe message pairs and finding the difference from the public $\sf Compare$ function.

We define the success probability of an adversary $\mathcal{A}$ in the IND game of the scheme $\Pi$ as
\[
	\Adv^{\sf IND}_{\Pi, \mathcal{A}^\mathcal{O}} := \Pr[{\sf IND}_{\Pi}(\mathcal{A}^\mathcal{O}) \to 1].
\]

The authentication scheme $\Pi$ is called \emph{IND} secure if for all PPT adversaries $\mathcal{A}$ when $\Pi$ has a non-trivial comparison secret key, or for all admissible PPT adversaries $\mathcal{A}$ when $\Pi$ does not have a non-trivial comparison secret key,
\[
	\Adv^{\sf IND}_{\Pi, \mathcal{A^\mathcal{O}}} = \negl.
\]



%-------------------
%% Reference List
\nocite{*}
\printbibliography


\end{document}
